%==========================================================
\chapter{جزئیات الگوریتم تولید نقشه اختلاف مکانی و تخمین عمق با بینایی استریو}
\label{appendix:stereo_depth_sgbm}

در این پیوست، فرایند کامل تولید داده‌های عمقی از تصاویر استریو تشریح می‌شود. هدف از این فرایند، تبدیل دو تصویر دوبعدی هم‌زمان، حاصل از دوربین چپ و راست، به نقشه اختلاف مکانی، نقشه عمق و در صورت نیاز ابرنقاط سه‌بعدی محیط است. این زنجیره شامل مراحل مدل‌سازی هندسی دوربین، کالیبراسیون، حذف اعوجاج، تصحیح استریو، محاسبه اختلاف مکانی، پالایش نقشه اختلاف مکانی و تبدیل اختلاف مکانی به عمق است.

%----------------------------------------------------------
\section{مدل هندسی دوربین و تصویر پرسپکتیو}

اساس محاسبه عمق در بینایی استریو بر «مدل دوربین سوراخ‌سوزنی%
\LTRfootnote{Pinhole Camera Model}» %
استوار است. در این مدل، یک نقطه سه‌بعدی در فضای دوربین با مختصات

\begin{equation}
	\mathbf{P}_c =
	\begin{bmatrix}
		X \\
		Y \\
		Z
	\end{bmatrix}
\end{equation}

روی صفحه تصویر به نقطه‌ای با مختصات پیکسلی \((u,v)\) نگاشت می‌شود. اگر فاصله کانونی دوربین در راستاهای افقی و عمودی به‌ترتیب \(f_x\) و \(f_y\)، و مرکز اپتیکی تصویر \((c_x,c_y)\) باشد، رابطه تصویرسازی به‌صورت زیر نوشته می‌شود:

\begin{equation}
	u = f_x \frac{X}{Z} + c_x,
	\qquad
	v = f_y \frac{Y}{Z} + c_y
	\label{eq:pinhole_projection_uv}
\end{equation}

این رابطه نشان می‌دهد که مختصات تصویر یک نقطه، به نسبت مختصات جانبی آن به عمق وابسته است. در فرم همگن، این نگاشت به‌صورت زیر بیان می‌شود:

\begin{equation}
	s
	\begin{bmatrix}
		u \\
		v \\
		1
	\end{bmatrix}
	=
	\mathbf{K}
	\begin{bmatrix}
		X \\
		Y \\
		Z
	\end{bmatrix},
	\label{eq:camera_projection_homogeneous}
\end{equation}

که در آن \(s\) ضریب مقیاس و \(\mathbf{K}\) ماتریس پارامترهای داخلی دوربین است:

\begin{equation}
	\mathbf{K} =
	\begin{bmatrix}
		f_x & 0 & c_x \\
		0 & f_y & c_y \\
		0 & 0 & 1
	\end{bmatrix}.
	\label{eq:camera_intrinsic_matrix}
\end{equation}

پارامترهای \(f_x\)، \(f_y\)، \(c_x\) و \(c_y\) از فرایند کالیبراسیون دوربین به‌دست می‌آیند. در محاسبه عمق مبتنی بر استریو، دانستن این پارامترها ضروری است؛ زیرا تبدیل اختلاف مکانی پیکسلی به عمق متریک مستقیماً به فاصله کانونی دوربین وابسته است.

%----------------------------------------------------------
\section{اعوجاج لنز و اصلاح تصویر}

در دوربین‌های واقعی، مدل سوراخ‌سوزنی به‌تنهایی کافی نیست، زیرا لنز دوربین باعث اعوجاج تصویر می‌شود. دو نوع رایج اعوجاج شامل «اعوجاج شعاعی%
\LTRfootnote{Radial Distortion}» %
و «اعوجاج مماسی%
\LTRfootnote{Tangential Distortion}» %
است. اگر مختصات نرمال‌شده یک نقطه در تصویر برابر با \((x,y)\) باشد، مقدار

\begin{equation}
	r^2 = x^2 + y^2
\end{equation}

تعریف می‌شود. در یک مدل متداول، اثر اعوجاج شعاعی و مماسی به‌صورت زیر نوشته می‌شود:

\begin{equation}
	x_d =
	x(1+k_1 r^2+k_2 r^4+k_3 r^6)
	+
	2p_1xy
	+
	p_2(r^2+2x^2)
	\label{eq:distorted_x}
\end{equation}

\begin{equation}
	y_d =
	y(1+k_1 r^2+k_2 r^4+k_3 r^6)
	+
	p_1(r^2+2y^2)
	+
	2p_2xy
	\label{eq:distorted_y}
\end{equation}

در این روابط، \(k_1\)، \(k_2\) و \(k_3\) ضرایب اعوجاج شعاعی و \(p_1\) و \(p_2\) ضرایب اعوجاج مماسی هستند. در مرحله کالیبراسیون، این ضرایب برای هر یک از دو دوربین چپ و راست تخمین زده می‌شوند. سپس تصاویر خام با استفاده از این ضرایب اصلاح می‌گردند تا نقاط تصویرشده تا حد امکان با مدل پرسپکتیو ایده‌آل سازگار شوند.

اصلاح اعوجاج برای محاسبه اختلاف مکانی اهمیت زیادی دارد؛ زیرا وجود اعوجاج باعث می‌شود نقاط متناظر در تصویر چپ و راست دقیقاً روی خطوط اپی‌پولار متناظر قرار نگیرند. در نتیجه، جستجوی متناظرها دشوارتر شده و نقشه اختلاف مکانی دچار خطا می‌شود.

%----------------------------------------------------------
\section{مدل هندسی سامانه استریو}

سامانه استریو از دو دوربین چپ و راست تشکیل شده است که با فاصله‌ای ثابت نسبت به یکدیگر نصب شده‌اند. این فاصله، فاصله مبنا%
\LTRfootnote{Baseline}
نامیده می‌شود و با \(B\) نمایش داده می‌شود. در حالت کلی، رابطه میان چارچوب دوربین چپ و راست با یک دوران \(\mathbf{R}\) و یک بردار انتقال \(\mathbf{T}\) توصیف می‌شود:

\begin{equation}
	\mathbf{P}_R =
	\mathbf{R}\mathbf{P}_L
	+
	\mathbf{T}
	\label{eq:left_right_extrinsic_relation}
\end{equation}

در این رابطه، \(\mathbf{P}_L\) مختصات نقطه نسبت به دوربین چپ و \(\mathbf{P}_R\) مختصات همان نقطه نسبت به دوربین راست است. بردار \(\mathbf{T}\) موقعیت نسبی مرکز دوربین راست نسبت به دوربین چپ را نشان می‌دهد. در یک آرایش استریوی ایده‌آل، دو دوربین کاملاً هم‌راستا هستند و انتقال تنها در راستای افقی اتفاق می‌افتد. در این حالت می‌توان نوشت:

\begin{equation}
	\mathbf{R} = \mathbf{I}
	\qquad
	\mathbf{T} =
	\begin{bmatrix}
		-B \\
		0 \\
		0
	\end{bmatrix}
\end{equation}

با این حال، در سامانه واقعی یا شبیه‌سازی‌شده، حتی اگر دوربین‌ها به‌صورت مکانیکی هم‌راستا نصب شوند، اختلافات کوچک در جهت‌گیری و موقعیت آن‌ها وجود دارد. بنابراین، قبل از محاسبه اختلاف مکانی، لازم است تصاویر چپ و راست وارد مرحله تصحیح استریو شوند.

%----------------------------------------------------------
\section{هندسه اپی‌پولار}

در بینایی استریو، هر نقطه سه‌بعدی توسط دو دوربین چپ و راست مشاهده می‌شود. اگر تصویر این نقطه در دوربین چپ برابر با \(\mathbf{x}_L\) و در دوربین راست برابر با \(\mathbf{x}_R\) باشد، این دو نقطه توسط قید اپی‌پولار به یکدیگر مرتبط می‌شوند:

\begin{equation}
	\mathbf{x}_R^T \mathbf{F} \mathbf{x}_L = 0
	\label{eq:epipolar_constraint}
\end{equation}

در این رابطه، \(\mathbf{F}\) ماتریس بنیادی%
\LTRfootnote{Fundamental Matrix}
است که رابطه هندسی بین دو تصویر را بیان می‌کند. اگر پارامترهای داخلی دوربین‌ها نیز معلوم باشند، رابطه اپی‌پولار در فضای نرمال‌شده با ماتریس ضروری%
\LTRfootnote{Essential Matrix}
بیان می‌شود:

\begin{equation}
	\mathbf{E} = [\mathbf{T}]_{\times}\mathbf{R}
	\label{eq:essential_matrix}
\end{equation}

که در آن \([\mathbf{T}]_{\times}\) ماتریس پادمتقارن متناظر با بردار انتقال \(\mathbf{T}\) است:

\begin{equation}
	[\mathbf{T}]_{\times} =
	\begin{bmatrix}
		0 & -T_z & T_y \\
		T_z & 0 & -T_x \\
		-T_y & T_x & 0
	\end{bmatrix}
\end{equation}

ماتریس بنیادی و ماتریس ضروری از رابطه زیر به یکدیگر مرتبط می‌شوند:

\begin{equation}
	\mathbf{F}
	=
	\mathbf{K}_R^{-T}
	\mathbf{E}
	\mathbf{K}_L^{-1}
	\label{eq:F_from_E}
\end{equation}

قید اپی‌پولار بیان می‌کند که اگر نقطه‌ای در تصویر چپ مشخص باشد، نقطه متناظر آن در تصویر راست روی یک خط مشخص، یعنی خط اپی‌پولار، قرار دارد. در تصاویر خام، این خط ممکن است مورب باشد؛ بنابراین برای هر پیکسل، جستجوی متناظر در تصویر دیگر باید در امتداد یک خط دوبعدی انجام شود. این کار از نظر محاسباتی پرهزینه و حساس به خطاست. هدف از تصحیح استریو آن است که خطوط اپی‌پولار را افقی کند تا جستجو فقط در راستای افقی تصویر انجام شود.

%----------------------------------------------------------
\section{تصحیح استریو و هم‌راستاسازی خطوط اپی‌پولار}

تصحیح استریو%
\LTRfootnote{Stereo Rectification}
فرایندی است که طی آن تصاویر چپ و راست به صفحات تصویری جدیدی نگاشت می‌شوند، به‌گونه‌ای که خطوط اپی‌پولار متناظر در هر دو تصویر افقی و هم‌تراز شوند. پس از این مرحله، اگر یک نقطه در تصویر چپ در مختصات \((u_L,v_L)\) قرار داشته باشد، نقطه متناظر آن در تصویر راست باید تقریباً روی همان سطر، یعنی در مختصات \((u_R,v_L)\)، قرار گیرد.

به بیان دیگر، پس از تصحیح استریو داریم:

\begin{equation}
	v_L \approx v_R
	\label{eq:rectified_same_row}
\end{equation}

بنابراین اختلاف مکانی فقط در راستای افقی تعریف می‌شود:

\begin{equation}
	d = u_L - u_R
	\label{eq:disparity_definition}
\end{equation}

در عمل، فرایند تصحیح استریو با استفاده از پارامترهای کالیبراسیون دو دوربین انجام می‌شود. این پارامترها شامل موارد زیر هستند:

\begin{equation}
	\mathbf{K}_L,\quad
	\mathbf{D}_L,\quad
	\mathbf{K}_R,\quad
	\mathbf{D}_R,\quad
	\mathbf{R},\quad
	\mathbf{T},
\end{equation}

که در آن \(\mathbf{K}_L\) و \(\mathbf{K}_R\) ماتریس‌های داخلی دوربین چپ و راست، \(\mathbf{D}_L\) و \(\mathbf{D}_R\) ضرایب اعوجاج، و \(\mathbf{R}\) و \(\mathbf{T}\) دوران و انتقال نسبی میان دو دوربین هستند.

خروجی مرحله تصحیح استریو معمولاً شامل ماتریس‌های دوران تصحیح‌شده \(\mathbf{R}_1\) و \(\mathbf{R}_2\)، ماتریس‌های تصویر جدید \(\mathbf{P}_1\) و \(\mathbf{P}_2\)، و ماتریس بازفرافکنی سه‌بعدی \(\mathbf{Q}\) است:

\begin{equation}
	\mathbf{R}_1,\quad
	\mathbf{R}_2,\quad
	\mathbf{P}_1,\quad
	\mathbf{P}_2,\quad
	\mathbf{Q}.
\end{equation}

پس از محاسبه این ماتریس‌ها، برای هر تصویر یک نگاشت پیکسلی ساخته می‌شود و تصویر خام به تصویر تصحیح‌شده تبدیل می‌گردد:

\begin{equation}
	I_L^{rect} = \mathcal{R}_L(I_L)
	\qquad
	I_R^{rect} = \mathcal{R}_R(I_R)
	\label{eq:rectified_images}
\end{equation}

در این رابطه، \(\mathcal{R}_L\) و \(\mathcal{R}_R\) نگاشت‌های تصحیح و حذف اعوجاج برای تصویر چپ و راست هستند.

%----------------------------------------------------------
\section{بهبود کنتراست تصاویر با روش \lr{CLAHE}}

پیش از محاسبه نقشه اختلاف مکانی، کیفیت ظاهری تصاویر ورودی نقش مهمی در پایداری فرایند تطابق دارد. الگوریتم‌های مبتنی بر تطابق استریو، از جمله \lr{SGBM}، برای یافتن نقاط متناظر میان تصویر چپ و راست به تغییرات شدت روشنایی، بافت تصویر و مرزهای محلی وابسته هستند. در نواحی کم‌کنتراست، یکنواخت یا دارای تغییرات نوری نامناسب، هزینه تطابق میان پیکسل‌های مختلف می‌تواند به یکدیگر نزدیک شود و در نتیجه مقدار اختلاف مکانی به‌درستی انتخاب نشود.

برای کاهش این مشکل، می‌توان از روش \lr{CLAHE} یا «برابرسازی هیستوگرام تطبیقی محدودشده%
\LTRfootnote{Contrast Limited Adaptive Histogram Equalization}»
استفاده کرد. این روش برخلاف برابرسازی هیستوگرام معمولی، به‌جای اعمال یک تبدیل سراسری روی کل تصویر، تصویر را به نواحی کوچک‌تر تقسیم کرده و کنتراست هر ناحیه را به‌صورت محلی افزایش می‌دهد. همچنین برای جلوگیری از تقویت بیش از حد نویز، مقدار افزایش کنتراست توسط یک آستانه محدود می‌شود.

در این پژوهش، اعمال \lr{CLAHE} پس از تصحیح استریو و پیش از اجرای الگوریتم \lr{SGBM} انجام می‌شود. به این ترتیب، تصاویر ورودی ابتدا از نظر هندسی هم‌راستا شده و سپس کنتراست محلی آن‌ها بهبود می‌یابد. اگر تصویر ورودی رنگی باشد، بهتر است \lr{CLAHE} روی مؤلفه روشنایی تصویر، مانند کانال خاکستری یا کانال (L) در فضای رنگی \lr{LAB}، اعمال شود تا ساختار رنگی تصویر دچار تغییر غیرضروری نشود.

به‌صورت کلی، اگر تصاویر تصحیح‌شده چپ و راست به‌ترتیب با ($I_L^{rect}$) و ($I_R^{rect}$) نمایش داده شوند، خروجی مرحله بهبود کنتراست به‌صورت زیر تعریف می‌شود:

\begin{equation}
	I_L^{clahe} = \mathcal{C}(I_L^{rect})
	\qquad
	I_R^{clahe} = \mathcal{C}(I_R^{rect})
	\label{eq:clahe_rectified_images}
\end{equation}

که در آن ($\mathcal{C}(\cdot)$) عملگر بهبود کنتراست محلی به روش \lr{CLAHE} است. این تصاویر بهبودیافته سپس به‌عنوان ورودی الگوریتم \lr{SGBM} مورد استفاده قرار می‌گیرند:

\begin{equation}
	D^{raw}
	=
	\mathcal{S}
	\left(
	I_L^{clahe},
	I_R^{clahe}
	\right)
	\label{eq:sgbm_after_clahe}
\end{equation}

در این رابطه، ($\mathcal{S}(\cdot)$) بیانگر فرایند محاسبه اختلاف مکانی با الگوریتم \lr{SGBM} و ($D^{raw}$) نقشه اختلاف مکانی خام است. استفاده از \lr{CLAHE} می‌تواند باعث افزایش پایداری تطابق در نواحی کم‌کنتراست شود، اما تنظیم نامناسب آن ممکن است نویز تصویر را نیز تقویت کند. بنابراین پارامترهایی مانند اندازه نواحی محلی و حد برش کنتراست باید به‌گونه‌ای انتخاب شوند که میان افزایش جزئیات تصویری و جلوگیری از تشدید نویز تعادل برقرار شود.

%----------------------------------------------------------
\section{اثبات رابطه عمق و اختلاف مکانی}

پس از تصحیح استریو، می‌توان رابطه عمق و اختلاف مکانی را از هندسه ساده مثلث‌های مشابه به‌دست آورد. فرض می‌شود دوربین چپ در مبدأ مختصات قرار دارد و دوربین راست به اندازه \(B\) در راستای محور \(x\) نسبت به آن جابه‌جا شده است. یک نقطه سه‌بعدی با مختصات

\begin{equation}
	\mathbf{P} =
	\begin{bmatrix}
		X \\
		Y \\
		Z
	\end{bmatrix}
\end{equation}

در تصویر چپ و راست مشاهده می‌شود. در تصویر چپ، مختصات افقی این نقطه برابر است با:

\begin{equation}
	u_L = f_x\frac{X}{Z} + c_x
	\label{eq:u_left_projection}
\end{equation}

از دید دوربین راست، مختصات افقی همان نقطه نسبت به مرکز دوربین راست برابر با \(X-B\) است. بنابراین، مختصات افقی تصویر آن در دوربین راست به‌صورت زیر نوشته می‌شود:

\begin{equation}
	u_R = f_x\frac{X-B}{Z} + c_x
	\label{eq:u_right_projection}
\end{equation}

اختلاف مکانی برابر اختلاف مختصات افقی نقطه متناظر در تصویر چپ و راست است:

\begin{equation}
	d = u_L - u_R
\end{equation}

با جایگذاری روابط \eqref{eq:u_left_projection} و \eqref{eq:u_right_projection} داریم:

\begin{equation}
	d =
	\left(
	f_x\frac{X}{Z} + c_x
	\right)
	-
	\left(
	f_x\frac{X-B}{Z} + c_x
	\right)
\end{equation}

با ساده‌سازی:

\begin{equation}
	d =
	f_x
	\left(
	\frac{X}{Z}
	-
	\frac{X-B}{Z}
	\right)
	=
	f_x
	\left(
	\frac{B}{Z}
	\right)
\end{equation}

در نتیجه:

\begin{equation}
	d = \frac{f_x B}{Z}
\end{equation}

و در نهایت رابطه عمق به‌صورت زیر به‌دست می‌آید:

\begin{equation}
	Z =
	\frac{f_x B}{d}
	\label{eq:depth_from_disparity_proof}
\end{equation}

این رابطه نشان می‌دهد که عمق با اختلاف مکانی رابطه معکوس دارد. هرچه جسم به دوربین نزدیک‌تر باشد، اختلاف مکانی آن در دو تصویر بیشتر است و هرچه جسم دورتر باشد، اختلاف مکانی کوچک‌تر می‌شود.

%----------------------------------------------------------
\section{بازسازی مختصات سه‌بعدی از نقشه اختلاف مکانی}

پس از محاسبه عمق \(Z\)، مختصات سه‌بعدی نقطه نیز از مدل دوربین به‌دست می‌آید. از رابطه تصویرسازی داریم:

\begin{equation}
	u = f_x\frac{X}{Z}+c_x
	\qquad
	v = f_y\frac{Y}{Z}+c_y
\end{equation}

بنابراین:

\begin{equation}
	X =
	\frac{(u-c_x)Z}{f_x}
	\qquad
	Y =
	\frac{(v-c_y)Z}{f_y}
	\label{eq:xyz_from_uvz}
\end{equation}

در نتیجه، هر پیکسل معتبر در نقشه اختلاف مکانی می‌تواند به یک نقطه سه‌بعدی تبدیل شود:

\begin{equation}
	(u,v,d)
	\rightarrow
	(X,Y,Z)
\end{equation}

با اعمال این تبدیل روی تمام پیکسل‌های معتبر تصویر، ابرنقاط سه‌بعدی محیط حاصل می‌شود. این ابرنقاط می‌تواند برای نمایش هندسه محیط، تخمین موقعیت موانع یا استخراج ویژگی‌های مکانی استفاده شود.

%----------------------------------------------------------
\section{ماتریس بازفرافکنی سه‌بعدی}

در بسیاری از پیاده‌سازی‌های بینایی ماشین، بازسازی سه‌بعدی از نقشه اختلاف مکانی با استفاده از ماتریس \(\mathbf{Q}\) انجام می‌شود. این ماتریس در مرحله تصحیح استریو تولید شده و رابطه میان مختصات پیکسلی و مختصات سه‌بعدی را به‌صورت همگن بیان می‌کند:

\begin{equation}
	\begin{bmatrix}
		X' \\
		Y' \\
		Z' \\
		W'
	\end{bmatrix}
	=
	\mathbf{Q}
	\begin{bmatrix}
		u \\
		v \\
		d \\
		1
	\end{bmatrix}
	\label{eq:q_reprojection}
\end{equation}

سپس مختصات سه‌بعدی واقعی از تقسیم بر مؤلفه همگن به‌دست می‌آید:

\begin{equation}
	X = \frac{X'}{W'}
	\qquad
	Y = \frac{Y'}{W'}
	\qquad
	Z = \frac{Z'}{W'}
	\label{eq:q_reprojection_normalized}
\end{equation}

در حالت ساده و ایده‌آل، این رابطه معادل همان روابط \eqref{eq:depth_from_disparity_proof} و \eqref{eq:xyz_from_uvz} است. مزیت استفاده از ماتریس \(\mathbf{Q}\) آن است که اثر پارامترهای کالیبراسیون، مرکز تصویر، فاصله مبنا و ماتریس‌های تصویر تصحیح‌شده را به‌صورت یکپارچه در نظر می‌گیرد.

%----------------------------------------------------------
\section{محاسبه نقشه اختلاف مکانی}

پس از تصحیح تصاویر، مسئله اصلی به یافتن نقطه متناظر برای هر پیکسل تصویر چپ در تصویر راست تبدیل می‌شود. از آنجا که تصاویر تصحیح شده‌اند، جستجو تنها در همان سطر انجام می‌شود. برای یک پیکسل \(p=(u,v)\) در تصویر چپ، هدف یافتن مقدار اختلاف مکانی \(d\) است به‌گونه‌ای که پیکسل متناظر آن در تصویر راست در موقعیت \((u-d,v)\) قرار گیرد.

به‌طور کلی، برای هر پیکسل و هر مقدار ممکن از اختلاف مکانی، یک هزینه تطابق تعریف می‌شود:

\begin{equation}
	C(p,d) =
	\mathcal{D}
	\left(
	I_L(u,v),
	I_R(u-d,v)
	\right)
	\label{eq:matching_cost_general}
\end{equation}

که در آن \(\mathcal{D}\) یک معیار اختلاف یا عدم شباهت میان شدت روشنایی یا ویژگی‌های تصویر است. ساده‌ترین معیارها شامل اختلاف قدرمطلق شدت روشنایی یا مجموع اختلاف قدرمطلق در یک پنجره محلی هستند:

\begin{equation}
	C_{\mathrm{SAD}}(p,d)
	=
	\sum_{(i,j)\in \Omega}
	\left|
	I_L(u+i,v+j)
	-
	I_R(u-d+i,v+j)
	\right|
	\label{eq:sad_cost}
\end{equation}

که در آن \(\Omega\) پنجره محلی اطراف پیکسل است. با این حال، استفاده مستقیم از چنین هزینه‌ای در بسیاری از نواحی تصویر، به‌ویژه در نواحی کم‌بافت، لبه‌ها، سایه‌ها و نقاط دارای انسداد، به خروجی ناپایدار منجر می‌شود. به همین دلیل، در بسیاری از کاربردهای عملی از روش‌های نیمه‌سراسری مانند \lr{SGBM} استفاده می‌شود.

%----------------------------------------------------------
\section{الگوریتم \lr{Semi-Global Block Matching}}

الگوریتم \lr{Semi-Global Block Matching} یا \lr{SGBM} یک روش نیمه‌سراسری برای محاسبه نقشه اختلاف مکانی است. ایده اصلی این روش، ترکیب هزینه تطابق محلی با قید هموارسازی اختلاف مکانی در چندین جهت تصویری است. این روش نسبت به تطابق بلوکی ساده پایدارتر است، اما هزینه محاسباتی آن نسبت به روش‌های سراسری کامل کمتر است.

در \lr{SGBM}، برای هر پیکسل \(p\) و هر مقدار اختلاف مکانی \(d\)، یک هزینه اولیه \(C(p,d)\) محاسبه می‌شود. سپس یک تابع انرژی برای کل نقشه اختلاف مکانی تعریف می‌شود:

\begin{equation}
	E(D)
	=
	\sum_{p}
	C(p,D_p)
	+
	\sum_{p}
	\sum_{q\in \mathcal{N}_p}
	\left[
	P_1 \, \mathbf{1}
	\left(
	|D_p-D_q|=1
	\right)
	+
	P_2 \, \mathbf{1}
	\left(
	|D_p-D_q|>1
	\right)
	\right]
	\label{eq:sgbm_energy}
\end{equation}

در این رابطه، \(D_p\) اختلاف مکانی انتخاب‌شده برای پیکسل \(p\)، و \(\mathcal{N}_p\) مجموعه همسایگان آن است. جمله اول، هزینه تطابق تصویر را نشان می‌دهد. جمله دوم، هموارسازی مکانی نقشه اختلاف مکانی را اعمال می‌کند. ضرایب \(P_1\) و \(P_2\) به‌ترتیب جریمه تغییرات کوچک و بزرگ در اختلاف مکانی هستند.

تفسیر این ضرایب به‌صورت زیر است:

\begin{itemize}
	\item \(P_1\): جریمه تغییرات کوچک اختلاف مکانی میان پیکسل‌های مجاور؛ این مقدار اجازه می‌دهد سطح‌های با شیب ملایم همچنان قابل مدل‌سازی باشند.
	\item \(P_2\): جریمه تغییرات بزرگ اختلاف مکانی؛ این مقدار از تغییرات ناگهانی غیرواقعی جلوگیری می‌کند، مگر در لبه‌های واقعی اجسام.
\end{itemize}

در روش سراسری کامل، کمینه‌سازی انرژی \eqref{eq:sgbm_energy} روی کل تصویر هزینه محاسباتی زیادی دارد. \lr{SGBM} به‌جای حل کامل دوبعدی، انرژی را در چندین مسیر یک‌بعدی در تصویر تقریب می‌زند.

%----------------------------------------------------------
\section{تجمیع هزینه در مسیرهای چندگانه}

در \lr{SGBM}، برای هر مسیر جهت‌دار \(r\)، هزینه تجمعی \(L_r(p,d)\) برای پیکسل \(p\) و اختلاف مکانی \(d\) به‌صورت بازگشتی محاسبه می‌شود:

\begin{equation}
	L_r(p,d)
	=
	C(p,d)
	+
	\min
	\begin{cases}
		L_r(p-r,d),\\
		L_r(p-r,d-1)+P_1,\\
		L_r(p-r,d+1)+P_1,\\
		\min\limits_{d'} L_r(p-r,d') + P_2
	\end{cases}
	-
	\min\limits_{d'} L_r(p-r,d')
	\label{eq:sgbm_path_cost}
\end{equation}

در این رابطه، \(p-r\) پیکسل قبلی در مسیر \(r\) است. جمله اول در عمل هزینه تطابق محلی را وارد می‌کند. جمله‌های داخل عملگر \(\min\) چهار حالت را در نظر می‌گیرند:

\begin{enumerate}
	\item اختلاف مکانی نسبت به پیکسل قبلی تغییر نکند؛
	\item اختلاف مکانی یک واحد کاهش یابد؛
	\item اختلاف مکانی یک واحد افزایش یابد؛
	\item اختلاف مکانی تغییر بزرگ‌تری داشته باشد.
\end{enumerate}

کم‌کردن مقدار \(\min_{d'}L_r(p-r,d')\) در انتهای رابطه برای جلوگیری از رشد عددی هزینه‌ها انجام می‌شود و مقدار بهینه اختلاف مکانی را تغییر نمی‌دهد.

پس از محاسبه هزینه تجمعی در چندین جهت، هزینه نهایی برای هر پیکسل و هر اختلاف مکانی از جمع مسیرها به‌دست می‌آید:

\begin{equation}
	S(p,d)
	=
	\sum_{r}
	L_r(p,d)
	\label{eq:sgbm_total_cost}
\end{equation}

در نهایت، اختلاف مکانی انتخاب‌شده برای پیکسل \(p\) برابر مقداری است که کمترین هزینه را دارد:

\begin{equation}
	D_p
	=
	\arg\min_{d}
	S(p,d)
	\label{eq:sgbm_winner_takes_all}
\end{equation}

این مرحله معمولاً با عنوان انتخاب برنده%
\LTRfootnote{Winner-Takes-All}
شناخته می‌شود.

%----------------------------------------------------------
\section{تنظیم پارامترهای مهم \lr{SGBM}}

کیفیت نقشه اختلاف مکانی به پارامترهای الگوریتم وابسته است. مهم‌ترین پارامترها عبارت‌اند از:

\begin{itemize}
	\item \(\mathrm{minDisparity}\): کمترین اختلاف مکانی قابل جستجو؛
	\item \(\mathrm{numDisparities}\): تعداد اختلاف مکانی‌های بررسی‌شده؛
	\item \(\mathrm{blockSize}\): اندازه پنجره محلی برای محاسبه هزینه تطابق؛
	\item \(P_1\): جریمه تغییرات کوچک اختلاف مکانی؛
	\item \(P_2\): جریمه تغییرات بزرگ اختلاف مکانی؛
	\item \(\mathrm{uniquenessRatio}\): آستانه یکتایی برای پذیرش پاسخ؛
	\item \(\mathrm{speckleWindowSize}\): اندازه پنجره حذف لکه‌های کوچک؛
	\item \(\mathrm{speckleRange}\): دامنه مجاز تغییر اختلاف مکانی در حذف لکه‌ها؛
	\item \(\mathrm{disp12MaxDiff}\): آستانه سازگاری چپ-راست.
\end{itemize}

به‌طور معمول، با افزایش \(\mathrm{blockSize}\)، نقشه اختلاف مکانی هموارتر اما جزئیات لبه‌ها ضعیف‌تر می‌شود. کاهش آن باعث حفظ جزئیات بیشتر می‌شود، اما حساسیت به نویز افزایش می‌یابد. همچنین مقدار بزرگ‌تر \(P_2\) باعث هموارتر شدن نقشه اختلاف مکانی می‌شود، اما اگر بیش از حد بزرگ انتخاب شود، لبه‌های واقعی اجسام از بین می‌روند.

در بسیاری از پیاده‌سازی‌ها، مقادیر \(P_1\) و \(P_2\) متناسب با تعداد کانال‌های تصویر و اندازه بلوک انتخاب می‌شوند:

\begin{equation}
	P_1 = 8 n_c b^2
	\qquad
	P_2 = 32 n_c b^2
	\label{eq:sgbm_p1_p2}
\end{equation}

که در آن \(n_c\) تعداد کانال‌های تصویر و \(b\) اندازه پنجره تطابق است. این روابط یک انتخاب رایج اولیه هستند و در عمل بسته به شرایط تصویر و نیاز سامانه تنظیم می‌شوند.

%----------------------------------------------------------
\section{اصلاح زیرپیکسلی اختلاف مکانی}

خروجی اولیه \lr{SGBM} معمولاً مقدار اختلاف مکانی را به‌صورت گسسته و در واحد پیکسل ارائه می‌کند. برای افزایش دقت عمق، می‌توان مقدار اختلاف مکانی را در اطراف کمینه هزینه به‌صورت زیرپیکسلی اصلاح کرد. اگر \(d_0\) اختلاف مکانی با کمترین هزینه باشد، و هزینه‌های متناظر با \(d_0-1\)، \(d_0\) و \(d_0+1\) به‌ترتیب \(S_{-}\)، \(S_0\) و \(S_{+}\) باشند، مقدار اصلاح‌شده می‌تواند با تقریب درجه دوم به‌صورت زیر محاسبه شود:

\begin{equation}
	d_{\mathrm{sub}}
	=
	d_0
	+
	\frac{
		S_{-}-S_{+}
	}{
		2(S_{-}-2S_0+S_{+})
	}
	\label{eq:subpixel_disparity}
\end{equation}

این اصلاح باعث می‌شود اختلاف مکانی از مقدار صحیح پیکسلی فراتر رفته و دقت عمق، به‌ویژه در فاصله‌های دورتر، بهبود یابد.

%----------------------------------------------------------
\section{سازگاری چپ-راست و حذف مقادیر نامعتبر}

یکی از مشکلات رایج در محاسبه اختلاف مکانی، نواحی انسداد%
\LTRfootnote{Occlusion}
است. در این نواحی، بخشی از صحنه در تصویر چپ دیده می‌شود اما در تصویر راست قابل مشاهده نیست، یا برعکس. در چنین شرایطی، متناظر واقعی برای برخی پیکسل‌ها وجود ندارد و خروجی تطابق می‌تواند اشتباه باشد.

برای کاهش این خطا، می‌توان نقشه اختلاف مکانی را یک‌بار از تصویر چپ به راست و یک‌بار از تصویر راست به چپ محاسبه کرد. سپس برای هر پیکسل، سازگاری دوطرفه بررسی می‌شود. اگر اختلاف مکانی چپ به راست برابر با \(d_L(u,v)\) باشد، نقطه متناظر در تصویر راست تقریباً در موقعیت \(u-d_L\) قرار می‌گیرد. اختلاف مکانی راست به چپ در آن نقطه باید با مقدار اولیه سازگار باشد:

\begin{equation}
	\left|
	d_L(u,v)
	-
	d_R(u-d_L(u,v),v)
	\right|
	<
	\epsilon_{LR}
	\label{eq:left_right_consistency}
\end{equation}

اگر این شرط برقرار نباشد، مقدار اختلاف مکانی آن پیکسل نامعتبر در نظر گرفته می‌شود. این آزمون به حذف خطاهای ناشی از انسداد، بازتاب و تطابق اشتباه کمک می‌کند.
%----------------------------------------------------------
\section{پالایش نقشه اختلاف مکانی با فیلتر \lr{WLS}}

نقشه اختلاف مکانی خام تولیدشده توسط \lr{SGBM} معمولاً شامل نویزهای موضعی، لکه‌های کوچک، نواحی نامعتبر و خطا در اطراف مرز اجسام است. اگر این نقشه به‌صورت مستقیم به عمق تبدیل شود، خطاهای کوچک در اختلاف مکانی، به‌ویژه در فاصله‌های دور، می‌توانند باعث خطای قابل توجه در تخمین فاصله شوند. به همین دلیل، پیش از تبدیل اختلاف مکانی به عمق، مرحله پالایش نقشه اختلاف مکانی انجام می‌شود.

یکی از روش‌های مؤثر برای این کار، فیلتر \lr{WLS} یا «کمینه مربعات وزن‌دار%
\LTRfootnote{Weighted Least Squares}»
است. هدف این فیلتر، هموارسازی نواحی نویزی نقشه اختلاف مکانی، در عین حفظ مرزهای واقعی اجسام است. برخلاف فیلترهای ساده‌ای مانند میانگین‌گیر یا فیلتر میانه، فیلتر \lr{WLS} از تصویر راهنما، معمولاً تصویر چپ تصحیح‌شده، برای تشخیص مرزهای تصویری استفاده می‌کند. بنابراین در نواحی یکنواخت، اختلاف مکانی هموار می‌شود، اما در نزدیکی لبه‌های مهم تصویر، از پخش‌شدن مقدار اختلاف مکانی میان دو جسم با عمق متفاوت جلوگیری می‌گردد.

اگر نقشه اختلاف مکانی خام با ($D^{raw}$) و نقشه اختلاف مکانی پالایش‌شده با ($D^{wls}$) نمایش داده شود، هدف فیلتر \lr{WLS} را می‌توان به‌صورت کمینه‌سازی تابع هزینه زیر بیان کرد:

\begin{equation}
	D^{wls}
	=
	\arg\min_D
	\left[
	\sum_{p}
	\left(
	D_p - D_p^{raw}
	\right)^2
	+
	\lambda
	\sum_{(p,q)\in \mathcal{N}}
	w_{pq}
	\left(
	D_p-D_q
	\right)^2
	\right]
	\label{eq:wls_disparity_energy}
\end{equation}

در این رابطه، جمله اول باعث می‌شود نقشه پالایش‌شده از مقدار خام اختلاف مکانی فاصله زیادی نگیرد. جمله دوم نقش هموارسازی مکانی را دارد. پارامتر ($\lambda$) شدت هموارسازی را کنترل می‌کند و ($w_{pq}$) وزن وابسته به شباهت پیکسل‌های مجاور است. این وزن معمولاً بر اساس اختلاف شدت روشنایی یا رنگ در تصویر راهنما تعریف می‌شود. به‌صورت مفهومی، برای دو پیکسل مجاور (p) و (q) می‌توان نوشت:

\begin{equation}
	w_{pq}
	=
	\exp
	\left(
	-
	\frac{
		\left|
		I_p-I_q
		\right|^2
	}{
		2\sigma_c^2
	}
	\right)
	\label{eq:wls_edge_weight}
\end{equation}

در این رابطه، ($I_p$) و ($I_q$) شدت روشنایی یا بردار رنگی پیکسل‌های متناظر در تصویر راهنما هستند و ($\sigma_c$) حساسیت فیلتر نسبت به تغییرات شدت تصویر را تعیین می‌کند. اگر دو پیکسل از نظر شدت روشنایی مشابه باشند، وزن ($w_{pq}$) بزرگ‌تر شده و هموارسازی میان آن‌ها بیشتر انجام می‌شود. اما اگر میان دو پیکسل یک لبه تصویری وجود داشته باشد، وزن کاهش می‌یابد و از هموارسازی شدید در دو طرف لبه جلوگیری می‌شود.

در پیاده‌سازی عملی، برای عملکرد بهتر فیلتر \lr{WLS} معمولاً علاوه بر نقشه اختلاف مکانی چپ به راست، نقشه اختلاف مکانی راست به چپ نیز محاسبه می‌شود. این کار امکان بررسی سازگاری چپ-راست و حذف بخشی از خطاهای ناشی از انسداد یا تطابق نادرست را فراهم می‌کند. سپس فیلتر \lr{WLS} با استفاده از تصویر چپ به‌عنوان تصویر راهنما، نقشه اختلاف مکانی نهایی را تولید می‌کند:

\begin{equation}
	D^{wls}
	=
	\mathcal{W}
	\left(
	D_L^{raw},
	D_R^{raw},
	I_L^{rect}
	\right)
	\label{eq:wls_filter_operation}
\end{equation}

که در آن ($D_L^{raw}$) نقشه اختلاف مکانی چپ به راست، ($D_R^{raw}$) نقشه اختلاف مکانی راست به چپ، ($I_L^{rect}$) تصویر راهنمای چپ، و ($\mathcal{W}(\cdot)$) عملگر فیلتر \lr{WLS} است.

نکته مهم آن است که فیلتر \lr{WLS} باید روی نقشه اختلاف مکانی اعمال شود، نه روی نقشه عمق. دلیل این موضوع آن است که رابطه عمق و اختلاف مکانی غیرخطی است و هموارسازی مستقیم عمق می‌تواند باعث ایجاد خطای جانبی در مرز اجسام شود. پس از پالایش اختلاف مکانی، نقشه عمق نهایی از رابطه زیر به‌دست می‌آید:

\begin{equation}
	Z(u,v)
	=
	\frac{f_xB}{D^{wls}(u,v)}
	\label{eq:depth_from_wls_disparity}
\end{equation}

بنابراین، فیلتر \lr{WLS} در زنجیره پردازش استریو به‌عنوان مرحله‌ای بین خروجی خام \lr{SGBM} و تبدیل اختلاف مکانی به عمق قرار می‌گیرد. این فیلتر با کاهش نویز و حفظ لبه‌های هندسی، باعث پایداری بیشتر تخمین فاصله هدف و موانع می‌شود.


%----------------------------------------------------------
\section{پالایش نقشه اختلاف مکانی}

خروجی خام \lr{SGBM} معمولاً شامل نویزهای نقطه‌ای، لکه‌های کوچک، نواحی نامعتبر و خطا در مرز اجسام است. به همین دلیل، پس از محاسبه نقشه اختلاف مکانی، معمولاً چند مرحله پالایش انجام می‌شود:

\begin{enumerate}
	\item حذف مقادیر نامعتبر یا منفی اختلاف مکانی؛
	\item حذف لکه‌های کوچک با استفاده از فیلتر \lr{speckle}؛
	\item اعمال فیلتر میانه برای حذف نویزهای نقطه‌ای؛
	\item استفاده از فیلترهای لبه‌نگهدار برای هموارسازی بدون تخریب مرز اجسام؛
	\item محدودسازی عمق به بازه معتبر مورد استفاده سامانه.
\end{enumerate}

در کاربردهای رباتیکی، این مرحله اهمیت زیادی دارد؛ زیرا خطای نقطه‌ای در نقشه اختلاف مکانی می‌تواند باعث تخمین نادرست فاصله مانع شود. با این حال، هموارسازی بیش از حد نیز مطلوب نیست، زیرا می‌تواند مرز موانع را جابه‌جا کرده و در تشخیص دقیق موقعیت آن‌ها خطا ایجاد کند. بنابراین، پالایش نقشه اختلاف مکانی باید تعادلی میان کاهش نویز و حفظ لبه‌های هندسی برقرار کند.

%----------------------------------------------------------
\section{تبدیل نقشه اختلاف مکانی به نقشه عمق}

پس از تولید و پالایش نقشه اختلاف مکانی \(D(u,v)\)، عمق هر پیکسل معتبر از رابطه زیر محاسبه می‌شود:

\begin{equation}
	Z(u,v)
	=
	\frac{f_x B}{D(u,v)}
	\label{eq:depth_map_from_disparity}
\end{equation}

در این رابطه، \(D(u,v)\) بر حسب پیکسل، \(f_x\) بر حسب پیکسل و \(B\) بر حسب متر است. بنابراین، خروجی \(Z(u,v)\) بر حسب متر خواهد بود. برای جلوگیری از تولید مقادیر نامعتبر، تنها پیکسل‌هایی استفاده می‌شوند که اختلاف مکانی آن‌ها مثبت و در بازه قابل قبول باشد:

\begin{equation}
	D(u,v) > 0
\end{equation}

از آنجا که رابطه عمق نسبت به اختلاف مکانی معکوس است، خطاهای کوچک در اختلاف مکانی در فاصله‌های دور اثر بزرگ‌تری بر خطای عمق می‌گذارند. اگر رابطه

\begin{equation}
	Z =
	\frac{f_x B}{d}
\end{equation}

را نسبت به \(d\) مشتق بگیریم، داریم:

\begin{equation}
	\frac{\partial Z}{\partial d}
	=
	-\frac{f_x B}{d^2}
	\label{eq:depth_derivative_disparity}
\end{equation}

با توجه به اینکه

\begin{equation}
	d =
	\frac{f_xB}{Z}
\end{equation}

می‌توان حساسیت عمق نسبت به خطای اختلاف مکانی را به‌صورت زیر نوشت:

\begin{equation}
	|\Delta Z|
	\approx
	\frac{Z^2}{f_xB}
	|\Delta d|
	\label{eq:depth_uncertainty}
\end{equation}

این رابطه نشان می‌دهد که عدم‌قطعیت عمق تقریباً با مربع فاصله افزایش می‌یابد. بنابراین، اندازه‌گیری عمق برای اجسام دورتر ذاتاً نویزی‌تر و کم‌دقت‌تر است. این نکته در طراحی سامانه ادراک و انتخاب محدوده معتبر عمق اهمیت دارد.

%----------------------------------------------------------
\section{تولید ابرنقاط سه‌بعدی}

پس از محاسبه نقشه عمق، مختصات سه‌بعدی هر پیکسل معتبر با استفاده از روابط تصویر معکوس محاسبه می‌شود:

\begin{equation}
	X(u,v)
	=
	\frac{(u-c_x)Z(u,v)}{f_x}
\end{equation}

\begin{equation}
	Y(u,v)
	=
	\frac{(v-c_y)Z(u,v)}{f_y}
\end{equation}

\begin{equation}
	Z(u,v)
	=
	\frac{f_xB}{D(u,v)}
\end{equation}

بنابراین هر پیکسل معتبر به یک نقطه سه‌بعدی تبدیل می‌شود:

\begin{equation}
	\mathbf{P}(u,v)
	=
	\begin{bmatrix}
		X(u,v) \\
		Y(u,v) \\
		Z(u,v)
	\end{bmatrix}
\end{equation}

مجموعه تمام این نقاط، ابرنقاط سه‌بعدی صحنه را تشکیل می‌دهد:

\begin{equation}
	\mathcal{P}
	=
	\left\{
	\mathbf{P}(u,v)
	\;|\;
	D(u,v)>0
	\right\}
\end{equation}

در سامانه‌های رباتیکی، معمولاً از تمام ابرنقاط برای تصمیم‌گیری مستقیم استفاده نمی‌شود، زیرا حجم داده زیاد است. در عوض، از نقشه عمق یا خلاصه‌ای از آن، مانند فاصله هدف، فاصله موانع یا موقعیت نسبی اجسام، به‌عنوان ورودی لایه ادراک و تصمیم‌گیری استفاده می‌شود.

%----------------------------------------------------------
\section{استخراج عمق هدف و موانع از نقشه اختلاف مکانی}

در کاربرد این پژوهش، هدف اصلی تولید نقشه عمق، استخراج فاصله و موقعیت نسبی هدف و موانع است. پس از آشکارسازی هدف یا مانع در تصویر، یک ناحیه دوبعدی متناظر با جعبه محدودکننده یا ناحیه هدف در تصویر به‌دست می‌آید. سپس مقادیر اختلاف مکانی داخل این ناحیه استخراج می‌شوند.

اگر مجموعه پیکسل‌های معتبر داخل ناحیه مورد نظر برابر با \(\Omega\) باشد، می‌توان مقدار اختلاف مکانی نماینده را به‌صورت میانگین یا میانه مقادیر معتبر محاسبه کرد:

\begin{equation}
	d_{\Omega}
	=
	\operatorname{median}
	\left(
	D(u,v)
	\;|\;
	(u,v)\in\Omega,\;
	D(u,v)>0
	\right)
	\label{eq:roi_disparity_median}
\end{equation}

یا در حالت میانگین:

\begin{equation}
	d_{\Omega}
	=
	\frac{1}{N_{\Omega}}
	\sum_{(u,v)\in\Omega}
	D(u,v)
	\label{eq:roi_disparity_mean}
\end{equation}

سپس عمق ناحیه از رابطه زیر محاسبه می‌شود:

\begin{equation}
	Z_{\Omega}
	=
	\frac{f_x B}{d_{\Omega}}
	\label{eq:roi_depth}
\end{equation}

استفاده از مقدار نماینده اختلاف مکانی، به‌جای استفاده مستقیم از تک‌پیکسل، باعث کاهش حساسیت به نویزهای نقطه‌ای و خطاهای موضعی نقشه اختلاف مکانی می‌شود. در نواحی دارای نویز زیاد، استفاده از میانه یا فیلترهای مقاوم نسبت به میانگین ساده پایدارتر است.

%----------------------------------------------------------
\section{خطاهای رایج در محاسبه عمق استریو}

اگرچه بینایی استریو یک روش مؤثر برای تخمین عمق است، اما خروجی آن تحت تأثیر چند منبع خطا قرار دارد:

\begin{enumerate}
	\item \textbf{نواحی کم‌بافت:} در سطوح یکنواخت، یافتن نقطه متناظر دشوار است.
	\item \textbf{لبه‌های عمقی:} در مرز اجسام، پنجره تطابق ممکن است شامل پیکسل‌هایی از دو عمق مختلف باشد.
	\item \textbf{انسداد:} برخی نقاط ممکن است فقط در یکی از دو تصویر دیده شوند.
	\item \textbf{بازتاب و نور متغیر:} تغییر روشنایی میان دو تصویر باعث افزایش خطای تطابق می‌شود.
	\item \textbf{کالیبراسیون نادرست:} خطا در \(f_x\)، \(B\)، یا تصحیح استریو مستقیماً باعث خطای عمق می‌شود.
	\item \textbf{خطای زیرپیکسلی در فاصله‌های دور:} طبق رابطه \eqref{eq:depth_uncertainty}، خطای کوچک در اختلاف مکانی برای اجسام دور به خطای عمق بزرگ‌تری تبدیل می‌شود.
\end{enumerate}

به همین دلیل، در سامانه‌های ناوبری، خروجی نقشه عمق معمولاً با فیلترهای مکانی و زمانی ترکیب می‌شود و برای تصمیم‌گیری نهایی از ویژگی‌های پالایش‌شده استفاده می‌گردد.

%----------------------------------------------------------
\section{روند اجرایی کامل تولید عمق}

روند کامل تولید عمق از تصاویر استریو را می‌توان به‌صورت الگوریتم زیر خلاصه کرد:

\begin{algorithm}[ht]
	\caption{تولید نقشه اختلاف مکانی و نقشه عمق از تصاویر استریو}
	\label{alg:stereo_depth_pipeline}
	\begin{algorithmic}[1]
		\REQUIRE تصویر چپ ($I_L$)، تصویر راست ($I_R$)، پارامترهای کالیبراسیون ($\mathbf{K}_L,\mathbf{D}_L,\mathbf{K}_R,\mathbf{D}_R,\mathbf{R},\mathbf{T}$)
		\ENSURE نقشه اختلاف مکانی پالایش‌شده ($D^{wls})$، نقشه عمق (Z)، و در صورت نیاز ابرنقاط ($\mathcal{P}$)
		
		
		\STATE حذف اعوجاج لنز از تصاویر چپ و راست
		\STATE محاسبه یا بارگذاری نگاشت‌های تصحیح استریو
		\STATE تولید تصاویر تصحیح‌شده \(I_L^{rect}\) و \(I_R^{rect}\)
		\STATE بهبود کنتراست محلی تصاویر تصحیح‌شده با روش \lr{CLAHE}
		\STATE محاسبه هزینه تطابق \(C(p,d)\) برای بازه مجاز اختلاف مکانی
		\STATE تجمیع هزینه‌ها در چندین مسیر با استفاده از الگوریتم \lr{SGBM}
		\STATE تولید نقشه اختلاف مکانی خام چپ به راست \(D_L^{raw}\)
		\STATE تولید نقشه اختلاف مکانی خام راست به چپ \(D_R^{raw}\) در صورت استفاده از آزمون سازگاری یا فیلتر \lr{WLS}
		\STATE اعمال آزمون سازگاری چپ-راست و حذف مقادیر نامعتبر
		\STATE پالایش نقشه اختلاف مکانی با فیلتر \lr{WLS} و تصویر چپ به‌عنوان تصویر راهنما
		\STATE حذف لکه‌های کوچک و مقادیر نامعتبر باقی‌مانده
		\STATE تبدیل اختلاف مکانی معتبر به عمق با رابطه \(Z=f_xB/D^{wls}\)
		\STATE بازسازی مختصات سه‌بعدی نقاط معتبر در صورت نیاز
		\STATE استخراج عمق هدف و موانع از نواحی آشکارسازی‌شده
		\RETURN \(D^{wls}\)، \(Z\)، \(\mathcal{P}\)
	\end{algorithmic}
	
\end{algorithm}
